\documentclass[10pt]{mypackage}

\usepackage[uprightgreeks,light]{kpfonts}
\renewcommand{\coloneq}{\coloneqq}
%\usepackage{homework}
\usepackage{notes}

\usepackage[ backend=bibtex, style = alphabetic, sorting=ynt ]{biblatex}
\addbibresource{all_references.bib}

\usepackage{parskip}

\fancyhf{}
\fancyhead[R]{Avinash Iyer}
\fancyhead[L]{Equivalence Relation von Neumann Algebras}
\fancyfoot[C]{\thepage}

\setcounter{secnumdepth}{0}

\begin{document}
\RaggedRight
\section{The Standard Representation and Tomita--Takesaki Theory}%
Before we may start understanding the von Neumann algebra of countable Borel equivalence relations, we discuss the standard representation of a finite von Neumann algebra $\left( M,\tau \right)$.
\begin{proposition}
  Let $M$ be a finite von Neumann algebra with trace $\tau$. Then, the map $x\mapsto x^{\ast}$ extends to an anti-linear isometry $J\colon L_2\left( M,\tau \right)\rightarrow L_2\left( M,\tau \right)$ with $JMJ = M'\cap B\left( L_2\left( M,\tau \right) \right)$. Here, $L_2\left( M,\tau \right)$ refers to the GNS representation of $M$ with respect to the faithful state $\tau$.
\end{proposition}
\begin{proof}
  For any $x\in M$, we have
  \begin{align*}
    \norm{x^{\ast}}_{\tau}^2 &= \tau\left( xx^{\ast} \right)\\
                             &= \tau\left( x^{\ast}x \right)\\
                             &= \norm{x}_{\tau}^2.
  \end{align*}
  Thus, this extends to an anti-linear isometry $J\colon L_2\left( M,\tau \right)\rightarrow L_2\left( M,\tau \right)$. We have $ \iprod{J\xi}{J\eta} = \iprod{\xi}{\eta} $ for all $\xi,\eta\in L_2\left( M,\tau \right)$.

  For any $\xi\in L_2\left( M,\tau \right)$, define the unbounded operators $L_{\xi}^{0}\colon M1_{\tau}\rightarrow L_2\left( M,\tau \right)$ and $R_{\xi}^{0}\colon M1_{\tau}\rightarrow L_2\left( M,\tau \right)$ by
  \begin{align*}
    L_{\xi}^0\left( x1_{\tau} \right) &= \left( Jx^{\ast}J \right)\xi\\
    R_{\xi}^{0}\left( x1_{\tau} \right) &= x\xi.
  \end{align*}
  Here, $1_{\tau}$ refers to the cyclic vector for the representation $\tau$.

  If $\left( x_i \right)_i\in M$ and $\eta\in L_2\left( M,\tau \right)$ are such that $\norm{x_i}_{\tau}\rightarrow 0$ and $L_{\xi}^0\left( x_i1_{\tau} \right)\rightarrow\eta$, then for all $z\in M$, we have
  \begin{align*}
    \left\vert \iprod{\eta}{z} \right\vert &= \lim_{i} \left\vert \iprod{\left( Jx_i^{\ast}J \right)\xi}{z} \right\vert\\
                                           &= \lim_{i} \left\vert \iprod{\xi}{zx_i^{\ast}} \right\vert\\
                                           &\leq \lim_{i} \norm{\xi}_{\tau}\norm{z}_{\op}\norm{x_i}_{\tau}\\
                                           &= 0,
  \end{align*}
  meaning that $L_{\xi}^0$ is closable. Completing the graph gives rise to a closed operator $L_{\xi}$, and similarly $R_{\xi}^0$ is closable with corresponding operator $R_{\xi}$. We have that $L_{x1_{\tau}} = x$ for all $x\in M$.

  Now, if $x,y\in M$ and $\xi\in L_2\left( M,\tau \right)$, then
  \begin{align*}
    \iprod{R_{J\xi}\left( x1_{\tau} \right)}{y1_{\tau}} &= \iprod{xJ\xi}{y1_{\tau}}\\
                                                        &= \iprod{y^{\ast}x1_{\tau}}{\xi}\\
                                                        &= \iprod{x1_{\tau}}{R_{\xi}\left( y1_{\tau} \right)},
  \end{align*}
  meaning that $R_{J\xi} = R_{\xi}^{\ast}$. Additionally,
  \begin{align*}
    \left( JL_{\xi}J \right)\left( x1_{\tau} \right) &= xJ\xi\\
                                                     &= R_{J\xi} \left( x1_{\tau} \right)\\
                                                     &= R_{\xi}^{\ast}\left( x1_{\tau} \right),
  \end{align*}
  so that $JL_{\xi}J = R_{\xi}^{\ast}$.

  If we set $ \widetilde{M} = \set{L_{\xi} | L_{\xi}\in B\left( L_2\left( M,\tau \right) \right)} $ and $ \widetilde{N} = \set{R_{\xi} | R_{\xi}\in B\left( L_2\left( M,\tau \right) \right)} $, then we have $J \widetilde{M}J = \widetilde{N}$.

  Therefore, it suffices to show that $M = \widetilde{M}$ and $M' = \widetilde{N}$. For this, we observe that $M\subseteq \widetilde{M}$ and $ \widetilde{N}\subseteq M' $. For any $x\in M'$, let $\xi = x1_{\tau}$. For any $y\in M$, we then have
  \begin{align*}
    R_{\xi}\left( y1_{\tau} \right) &= y\xi\\
                                    &= yx1_{\tau}\\
                                    &= xy1_{\tau},
  \end{align*}
  meaning that $R_{\xi} = x$, so $M'\subseteq \widetilde{N}$. For any $L_{\xi}\in \widetilde{M}$, we then have $L_{\xi}\in \widetilde{N}' = M'' = M$, so $\widetilde{M}\subseteq M$.
\end{proof}
The standard representation in fact generalizes to a wider setting --- namely, $\tau$ can be merely a faithful normal state instead of having to be a trace. This is the realm of Tomita--Takesaki Theory, and we scratch the surface of it here with a discussion of Hilbert algebras.
\section{Equivalence Relation von Neumann Algebras}%
We start by considering a type of ``mass transport principle'' on an equivalence relation $\mathcal{R}\subseteq X\times X$, where $\left( X,\mu \right)$ is some probability space. From now on, we will disregard null sets (i.e., ``every'' means ``almost every''). We will also assume that (almost) every orbit of $\mathcal{R}$ is infinite.
\begin{definition}
  If $\mathcal{R}\subseteq X\times X$ is an equivalence relation, and $\mathcal{W}\subseteq \mathcal{R}$ is Borel, then we define
  \begin{align*}
    \nu\left( \mathcal{W} \right) &= \int_{X}^{} \mu\left( \set{y\in X | \left( x,y \right)\in \mathcal{W}} \right)\:d\mu(x).
  \end{align*}
\end{definition}
Note that since $\mu$ is invariant under $\mathcal{R}$, we have that
\begin{align*}
  \int_{X}^{} \mu\left( \set{y\in X | \left( x,y \right)\in \mathcal{W}} \right)\:d\mu(x) &= \int_{X}^{} \mu\left( \set{x\in X | \left( x,y \right)\in \mathcal{W}} \right)\:d\mu(y).
\end{align*}
Note that if $f\colon X_1\rightarrow X_2$ is a Borel isomorphism, then there is a unique $\pi_{f}\colon L_{\infty}\left( X_1 \right)\rightarrow L_{\infty}\left( X_2 \right)$ given by $\pi_{f}\left( \phi \right)\left( x \right) = \phi\left( f^{-1}(x) \right)$.

Now, consider the space $L_2\left( \mathcal{R},\nu \right)$. If $\phi\in \left[ \left[ \mathcal{R} \right] \right]$, and $\xi\in L_2\left( \mathcal{R},\nu \right)$, define
\begin{align*}
  \left( u_{\phi}\xi \right)\left( x,y \right) &= \chi_{\operatorname{ran}\left( \phi \right)}(x)\xi\left( \phi^{-1}\left( x \right),y \right).
\end{align*}
That is,
\begin{align*}
  \left( u_{\phi}\xi \right)\left( x,y \right) &= \xi\left( \phi^{-1}(x),y \right)
\end{align*}
if $x\in \operatorname{ran}\left( \phi \right)$ and $0$ otherwise. The von Neumann algebra $L_{\infty}(X)$ acts on $L_2\left( \mathcal{R},\nu \right)$ by multiplication in the $x$ coordinate:
\begin{align*}
  \left( F\xi \right)\left( x,y \right) &= F\left( x \right)\xi\left( x,y \right).
\end{align*}
If $\phi\in \left[ \left[ \mathcal{R} \right] \right]$ and $F\in L_{\infty}\left( X \right)$, then $F_{\phi}\in L_{\infty}\left( X \right)$ can be defined by
\begin{align*}
  F_{\phi}(x) &= \chi_{\operatorname{ran}\left( \phi \right)}(x)F\left( \phi^{-1}\left( x \right) \right).
\end{align*}
That is, $F_{\phi}(x) = F\left( \phi^{-1}\left( x \right) \right)$ if $x\in \operatorname{ran}\left( \phi \right)$ and $0$ otherwise. We call this the left-regular representation of $\mathcal{R}$.

Similarly, it is possible to define the right-regular representation of $\mathcal{R}$ by taking $\left( v_{\phi}\xi \right)\left( x,y \right) = \1_{\operatorname{ran}\left( \phi \right)}(y)\xi\left( x,\phi^{-1}\left( y \right) \right)$.
\begin{proposition}
  If $\phi,\psi\in \left[ \left[ R \right] \right]$, the following hold:
  \begin{enumerate}[(i)]
    \item $u_{\phi}u_{\psi} = u_{\phi\psi}$;
    \item $u_{\phi}^{\ast} = u_{\phi^{-1}}$;
    \item $u_{\phi}^{\ast}u_{\psi} = \chi_{\operatorname{dom}\left( \phi \right)}$ and $u_{\phi}u_{\phi}^{\ast} = \chi_{\operatorname{ran}\left( \phi \right)}$;
    \item $u_{\phi}F = F_{\phi}u_{\phi}$ for any $F\in L_{\infty}\left( X \right)$.
  \end{enumerate}
\end{proposition}
\begin{definition}
  Define the von Neumann algebra $L\left( \mathcal{R} \right)\subseteq B\left( L_2\left( \mathcal{R},\nu \right) \right)$ to be the von Neumann algebra generated by the set $\set{u_{\phi} | \phi\in \left[ \left[ \mathcal{R} \right] \right]}$.

  Likewise, define $R\left( \mathcal{R} \right)\subseteq B\left( L_2\left( \mathcal{R},\nu \right) \right)$ to be the von Neumann algebra generated by $\set{v_{\phi} | \phi\in \left[ \left[ \mathcal{R} \right] \right]}$.
\end{definition}
We will specialize this for the case of the von Neumann algebra $L\left( \mathcal{R} \right)$. Observe that $L\left( \mathcal{R} \right)\subseteq R\left( \mathcal{R} \right)'$, so we may define the canonical anti-unitary $J\colon L_2\left( \mathcal{R},\nu \right)\rightarrow L_2\left( \mathcal{R},\nu \right)$ by $\left( J\xi \right)\left( x,y \right) = \overline{\xi\left( y,x \right)}$.
\begin{proposition}
  Let $\xi_0 = \chi_{\Delta}\in L_2\left( \mathcal{R},\nu \right)$, where $\Delta\subseteq \mathcal{R}$ is the diagonal.
  \begin{enumerate}[(i)]
    \item We have that $\xi_0$ is a cyclic and separating vector for $L\left( \mathcal{R} \right)$.
    \item The vector state $ \tau = \iprod{\cdot \xi_0}{\xi_0} $ is a faithful normal trace on $L\left( \mathcal{R} \right)$. In particular, $L\left( \mathcal{R} \right)$ is a finite von Neumann algebra.
    \item For any $\phi\in \left[ \left[ \mathcal{R} \right] \right]$, we have $Ju_{\phi}J = v_{\phi}$, and in particular, $L\left( \mathcal{R} \right) = R\left( \mathcal{R} \right)'$.
  \end{enumerate}
\end{proposition}
\begin{proof}\hfill
  \begin{enumerate}[(i)]
    \item Let $\xi_0 = \chi_{\Delta}$. For any $\phi\in \left[ \left[ \mathcal{R} \right] \right]$, we have
      \begin{align*}
        \chi_{\operatorname{graph}\left( \phi \right)} &= u_{\phi^{-1}}\xi_0.
      \end{align*}
      Since $\mathcal{R}$ can be written as a countable disjoint union of the graphs of Borel isomorphisms, $\mathcal{R} = \bigsqcup_{n}\operatorname{graph}\left( g_n \right)$, so we may write any $W\subseteq \mathcal{R}$ as $W = \bigsqcup_{n}W_n$ with $W_n = W\cap \operatorname{graph}\left( g_n \right)$. Since $W_n$ is the graph of a partial Borel isomorphism, we have that $\chi_W\in \overline{L\left( \mathcal{R} \right)\xi_0}$, so $\xi_0$ is cyclic for $L\left( \mathcal{R} \right)$.
    \item It suffices to show that $\tau\left( u_{\phi}u_{\psi} \right) = \tau\left( u_{\psi}u_{\phi} \right)$ for all $\phi,\psi\in \left[ \left[ \mathcal{R} \right] \right]$. We have
      \begin{align*}
        \tau\left( u_{\phi}u_{\psi} \right) &= \mu\left( \set{x\in X | \left( \phi\circ\psi \right)^{-1}(x) = x} \right)\\
                                            &= \mu\left( \set{x\in X | \psi^{-1}\circ\phi^{-1}(x) = x} \right)\\
                                            &= \mu\left( \set{x\in X | \phi^{-1}\circ\psi^{-1}\left( \phi^{-1}\left( x \right) \right) = \phi^{-1}\left( x \right)} \right)\\
                                            &= \mu\left( \set{x\in X | \left( \psi\circ\phi \right)^{-1}\left( x \right) = x} \right)\\
                                            &= \tau\left( u_{\psi}u_{\phi} \right).
      \end{align*}
  \end{enumerate}
\end{proof}
Let $A = L_{\infty}\left( X \right)$ and $M = L\left( \mathcal{R} \right)$. We have that $A\subseteq M$. There must then exist a unique $\tau$-preserving normal conditional expectation $E_A\colon M\rightarrow A$, by some results in \href{https://ai.avinash-iyer.com/Classes_and_Homework/After\%20College/Other\%20Notes/vna_projections.pdf}{von Neumann algebra theory}. It suffices to compute $E_A\left( u_{\phi} \right)$ for any $\phi\in \left[ \left[ \mathcal{R} \right] \right]$. Denote by $e_D\colon L_2\left( \mathcal{R} \right)\rightarrow L_2\left( \Delta, \right)$ the orthogonal projection onto the subspace of the diagonal.
\nocite{takesaki_1,takesaki_2,von_neumann_algebras_peterson,ii1_factors_houdayer}
\printbibliography
\end{document}
